%=============================================================
%  PEDAGOGICAL SUPPORT MATERIAL
%  Quantum Mechanics & Quantum Chemistry Programme
%  Learning Tiers + Assessment Artifact Set
%=============================================================
\documentclass[11pt,a4paper]{article}

%--- Packages ------------------------------------------------
\usepackage[T1]{fontenc}
\usepackage[utf8]{inputenc}
\usepackage[margin=2.5cm]{geometry}
\usepackage{amsmath,amssymb}
\usepackage{booktabs}
\usepackage{array}
\usepackage{longtable}
\usepackage{tabularx}
\usepackage{multirow}
\usepackage{enumitem}
\usepackage{sectsty}
\usepackage{fancyhdr}
\usepackage[dvipsnames,svgnames,table]{xcolor}
\usepackage{hyperref}
\usepackage{parskip}
\usepackage{tcolorbox}
\usepackage{mdframed}
\tcbuselibrary{skins,breakable,listings}

%--- Colours -------------------------------------------------
\definecolor{darkblue}{RGB}{10,40,90}
\definecolor{midblue}{RGB}{30,80,160}
\definecolor{lightblue}{RGB}{220,233,255}
\definecolor{accentgold}{RGB}{180,140,20}
\definecolor{lightgold}{RGB}{255,248,220}
\definecolor{darkgray}{RGB}{60,60,60}
\definecolor{lightgray}{RGB}{245,245,245}
\definecolor{tier1col}{RGB}{180,230,180}
\definecolor{tier2col}{RGB}{150,210,240}
\definecolor{tier3col}{RGB}{200,180,240}
\definecolor{tier4col}{RGB}{255,210,150}
\definecolor{tier5col}{RGB}{255,170,150}
\definecolor{artteal}{RGB}{20,110,100}
\definecolor{lightartteal}{RGB}{210,240,238}
\definecolor{artpurple}{RGB}{100,30,130}
\definecolor{lightartpurple}{RGB}{235,215,245}

%--- Section fonts -------------------------------------------
\sectionfont{\Large\bfseries\color{darkblue}\sectionrule{0ex}{0pt}{-1ex}{1.5pt}}
\subsectionfont{\large\bfseries\color{midblue}}
\subsubsectionfont{\normalsize\bfseries\color{darkgray}}

%--- Hyperref ------------------------------------------------
\hypersetup{
  colorlinks=true,
  linkcolor=midblue,
  urlcolor=midblue,
  citecolor=midblue,
  pdftitle={Pedagogical Support Material -- QM Programme},
}

%--- Header / Footer -----------------------------------------
\pagestyle{fancy}
\fancyhf{}
\fancyhead[L]{\small\color{darkgray}\textit{QM Programme --- Pedagogical Support}}
\fancyhead[R]{\small\color{darkgray}\thepage}
\fancyfoot[C]{\small\color{darkgray}\textit{For instructors and students}}
\renewcommand{\headrulewidth}{0.4pt}

%--- tcolorbox styles ----------------------------------------
\tcbset{
  tierbox/.style={
    enhanced, breakable,
    colframe=#1!80!black,
    colback=#1!15!white,
    fonttitle=\bfseries\color{white},
    coltitle=white,
    attach boxed title to top left={yshift=-2mm, xshift=4mm},
    boxed title style={colback=#1!80!black, sharp corners, boxrule=0.5pt},
    sharp corners=south, boxrule=0.8pt,
    left=6pt, right=6pt, top=8pt, bottom=6pt,
  },
  artifactbox/.style={
    enhanced, breakable,
    colback=lightartteal, colframe=artteal,
    fonttitle=\bfseries\color{white}, coltitle=white,
    attach boxed title to top left={yshift=-2mm, xshift=4mm},
    boxed title style={colback=artteal, colframe=artteal!80!black,
      sharp corners, boxrule=0.5pt},
    sharp corners=south, boxrule=0.8pt,
    left=6pt, right=6pt, top=8pt, bottom=6pt,
  },
  rubricbox/.style={
    enhanced,
    colback=lightartpurple, colframe=artpurple,
    fonttitle=\bfseries\color{white}, coltitle=white,
    attach boxed title to top left={yshift=-2mm, xshift=4mm},
    boxed title style={colback=artpurple, sharp corners, boxrule=0.5pt},
    sharp corners=south, boxrule=0.7pt,
    left=5pt, right=5pt, top=6pt, bottom=5pt,
  },
  goldbox/.style={
    enhanced, breakable,
    colback=lightgold, colframe=accentgold,
    boxrule=0.8pt,
    left=6pt, right=6pt, top=6pt, bottom=6pt,
    fontupper=\small,
  },
  notebox/.style={
    enhanced,
    colback=lightblue, colframe=midblue,
    boxrule=0.6pt, rounded corners,
    left=5pt, right=5pt, top=4pt, bottom=4pt,
    fontupper=\small\itshape,
  },
}

\newcommand{\instructornote}[1]{%
  \begin{tcolorbox}[notebox]
    \textbf{Instructor note:} #1
  \end{tcolorbox}%
}

\newcommand{\capabilityrow}[2]{\textbf{#1} & #2 \\}

%--- Tier colour commands ------------------------------------
\newcommand{\tiertag}[2]{%
  \colorbox{#1}{\small\bfseries\color{darkblue}\ #2\ }%
}

%=============================================================
\begin{document}
%=============================================================

%--- Title Page ----------------------------------------------
\begin{titlepage}
\pagecolor{darkblue}
\color{white}
\vspace*{2.5cm}
\begin{center}
  {\fontsize{26}{32}\selectfont\bfseries Pedagogical Support Material}\\[0.5cm]
  {\fontsize{16}{22}\selectfont Quantum Mechanics \& Quantum Chemistry Programme}\\[0.6cm]
  {\color{accentgold}\rule{10cm}{2pt}}\\[0.6cm]
  {\large\itshape For Students and Instructors}\\[2.5cm]
  \begin{tabular}{cl}
    \textbf{Section~1:} & Learning Tiers (HS $\to$ PhD) --- Definitions, Competencies\\[4pt]
                        & \quad and a Worked Example from Bra-Ket Module I.2\\[8pt]
    \textbf{Section~2:} & Assessment Artifact Set --- Full Taxonomy,\\[4pt]
                        & \quad Templates, Rubrics, and Packaging Standard\\
  \end{tabular}\\[3cm]
  {\color{accentgold}\rule{10cm}{1pt}}\\[0.5cm]
  {\small\itshape This document is a companion to the Course Syllabus.\\
  It sets out the shared language for \emph{how} material is taught and assessed\\
  across all Series and Modules of the programme.}
\end{center}
\vfill
\end{titlepage}
\nopagecolor

%--- Table of Contents ---------------------------------------
\tableofcontents
\newpage

%=============================================================
\section*{Introduction}
\addcontentsline{toc}{section}{Introduction}
%=============================================================

\begin{tcolorbox}[goldbox]
This document serves two audiences simultaneously.

\medskip
\noindent\textbf{For students:} It explains exactly what is expected of you at each stage of
your education, and what kinds of tasks you will be asked to perform. Understanding the
tier system and the artifact taxonomy before a course begins removes ambiguity about
difficulty levels and grading expectations.

\medskip
\noindent\textbf{For instructors:} It provides a reusable, course-independent specification
for how to differentiate content and assessment. The tier system and artifact set described
here apply uniformly across \emph{every lecture} in \emph{every module} of the programme.
Consistency of language across Series~I, II, and III is a deliberate design choice.
\end{tcolorbox}

\medskip
\noindent\textbf{How the two sections connect.}
Section~1 defines the \emph{audience} dimension: who is in the room and what they can be
expected to bring. Section~2 defines the \emph{assessment} dimension: what tasks reveal
mastery at each level.
Together they answer the two questions every educator and student needs:
\begin{itemize}[leftmargin=1.8em]
  \item \textit{What level is this material pitched at?}
  \item \textit{What does ``mastery'' look like for me at my level?}
\end{itemize}

\newpage
%#############################################################
\section{Learning Tiers}
\label{sec:tiers}
%#############################################################

\begin{tcolorbox}[goldbox]
\textbf{Overview.}
The programme is designed to be taught to a \emph{mixed cohort}: a single lecture or module
may be attended simultaneously by high-school students with strong mathematical aptitude
and PhD researchers extending their foundations.
The five-tier system ensures every student finds appropriate entry points, stretch goals,
and assessment tasks, without fragmenting the cohort into separate tracks.

\medskip
\noindent The tiers are \emph{not grade labels}. They describe the prerequisite
competency set a student brings and therefore what kind of engagement they can
immediately sustain. A student may operate at different tiers for different topics
(e.g.\ Tier~3 for linear algebra, Tier~2 for group theory).
\end{tcolorbox}

%-------------------------------------------------------------
\subsection{Tier Ladder at a Glance}

\begin{center}
\begin{tabular}{clp{8.5cm}}
\toprule
\textbf{Tier} & \textbf{Label} & \textbf{One-line description}\\
\midrule
1 & High School (HS)    & Amplitudes are complex arrows; probabilities are squared lengths\\
2 & Beginning UG (BegUG) & Bases and components; compute probabilities with $\langle n|\psi\rangle$\\
3 & Advanced UG (AdvUG)  & Inequalities; geometry and internal consistency of the formalism\\
4 & Master's (MSc)       & Rays and projective space; symmetries via Wigner's theorem\\
5 & PhD                  & Rigorous Hilbert space; separability; domain subtleties\\
\bottomrule
\end{tabular}
\end{center}

\medskip
The five tiers form a \emph{cumulative} ladder: every competency at Tier $k$ is assumed
at Tier $k{+}1$.
Instructors should always state which tier(s) a given lecture segment targets.

%-------------------------------------------------------------
\subsection{Tier Definitions and Competencies}

%--- Tier 1 --------------------------------------------------
\begin{tcolorbox}[tierbox=tier1col, title={Tier 1 --- High School (HS)}]
\textbf{Goal:} Build the \emph{intuition tools} needed for amplitudes and interference.

\medskip
\noindent\textbf{Core competencies:}
\begin{itemize}[leftmargin=1.5em,itemsep=2pt]
  \item Complex numbers as ``2D arrows'': $a+ib$, magnitude $|z|$, phase $\arg z$
  \item \textbf{Magnitude squared} as probability: $P \propto |A|^2$
  \item \textbf{Vector addition analogy}: adding complex amplitudes before squaring gives
        interference
  \item Basic probability: events, normalisation (``total probability = 1'')
\end{itemize}

\medskip
\noindent\textbf{You should be able to:}
\begin{itemize}[leftmargin=1.5em,itemsep=2pt]
  \item Compute $|a+ib|^2$
  \item Understand why $|A_1+A_2|^2\neq |A_1|^2+|A_2|^2$ in general
\end{itemize}
\end{tcolorbox}

\medskip
%--- Tier 2 --------------------------------------------------
\begin{tcolorbox}[tierbox=tier2col, title={Tier 2 --- Beginning Undergraduate (BegUG)}]
\textbf{Goal:} Become fluent in the \emph{working language} of bras, kets, bases,
and probabilities.

\medskip
\noindent\textbf{Core competencies:}
\begin{itemize}[leftmargin=1.5em,itemsep=2pt]
  \item Normalisation: $\langle\psi|\psi\rangle=1$
  \item Orthonormal bases: $\langle m|n\rangle=\delta_{mn}$
  \item Components as inner products: $c_n=\langle n|\psi\rangle$
  \item Probabilities: $P(n)=|c_n|^2$
  \item ``Insert identity'' in discrete form: $\displaystyle\sum_n |n\rangle\langle n|=\hat{\mathbb{1}}$
\end{itemize}

\medskip
\noindent\textbf{You should be able to:}
\begin{itemize}[leftmargin=1.5em,itemsep=2pt]
  \item Expand a state in a basis and compute probabilities
  \item Translate between $|\psi\rangle$ and its coordinate column vector in a chosen basis
\end{itemize}
\end{tcolorbox}

\medskip
%--- Tier 3 --------------------------------------------------
\begin{tcolorbox}[tierbox=tier3col, title={Tier 3 --- Advanced Undergraduate (AdvUG)}]
\textbf{Goal:} Understand why the formalism is \emph{internally consistent} and geometric.

\medskip
\noindent\textbf{Core competencies:}
\begin{itemize}[leftmargin=1.5em,itemsep=2pt]
  \item \textbf{Cauchy--Schwarz inequality}:
        $|\langle\phi|\psi\rangle|^2 \le \langle\phi|\phi\rangle\langle\psi|\psi\rangle$
  \item \textbf{Triangle inequality} for norms:
        $\|\phi+\psi\|\le\|\phi\|+\|\psi\|$
  \item Metric/geometric viewpoint:
        overlaps behave like cosines of angles (up to phase);
        distance from $\|\phi-\psi\|$ (with phase caveats)
\end{itemize}

\medskip
\noindent\textbf{You should be able to:}
\begin{itemize}[leftmargin=1.5em,itemsep=2pt]
  \item Prove basic inequalities in Hilbert space
  \item Explain why overlaps and probabilities are bounded by 1 for normalised states
  \item See superposition as linear structure, not mystery
\end{itemize}
\end{tcolorbox}

\medskip
%--- Tier 4 --------------------------------------------------
\begin{tcolorbox}[tierbox=tier4col, title={Tier 4 --- Master's (MSc)}]
\textbf{Goal:} Shift from ``vectors as states'' to \emph{rays as states} and connect
to symmetry theory.

\medskip
\noindent\textbf{Core competencies:}
\begin{itemize}[leftmargin=1.5em,itemsep=2pt]
  \item \textbf{Rays / projective Hilbert space}: physical states are equivalence classes
        $|\psi\rangle\sim e^{i\theta}|\psi\rangle$; pure states are points in
        $\mathbb{P}(\mathcal{H})$
  \item \textbf{Wigner theorem preview}: transformations preserving transition probabilities
        $|\langle\phi|\psi\rangle|^2$ are implemented by \emph{unitary or antiunitary} operators
\end{itemize}

\medskip
\noindent\textbf{You should be able to:}
\begin{itemize}[leftmargin=1.5em,itemsep=2pt]
  \item Explain why global phase is unphysical and why rays are the correct objects
  \item State what Wigner's theorem says and why it matters for symmetries
\end{itemize}
\end{tcolorbox}

\medskip
%--- Tier 5 --------------------------------------------------
\begin{tcolorbox}[tierbox=tier5col, title={Tier 5 --- PhD}]
\textbf{Goal:} Make all foundational statements \emph{mathematically precise} and know
where physics shorthand hides subtleties.

\medskip
\noindent\textbf{Core competencies:}
\begin{itemize}[leftmargin=1.5em,itemsep=2pt]
  \item Precise definition of \textbf{Hilbert space} (inner product + completeness)
  \item \textbf{Separability} (existence of a countable orthonormal basis / countable dense subset)
        and why it is assumed in most QM
  \item Rays as \textbf{1D subspaces} (not just ``vectors up to phase'')
  \item Operator subtleties: Hermitian vs.\ self-adjoint (domains matter);
        generalised eigenvectors and distributional states (rigged Hilbert spaces)
\end{itemize}

\medskip
\noindent\textbf{You should be able to:}
\begin{itemize}[leftmargin=1.5em,itemsep=2pt]
  \item Distinguish informal bra-ket manipulations from rigorous operator statements
  \item Discuss why separability matters for spectral decompositions
  \item Frame state space as a projective object and symmetry as structure on that space
\end{itemize}
\end{tcolorbox}

%-------------------------------------------------------------
\subsection{What Changes Across Tiers (Without Changing the Artifact)}

\instructornote{%
  The same assignment can be given to all tiers simultaneously. What changes is the
  \emph{language}, the required \emph{rigor}, and the \emph{scope}.
  The table below illustrates this for the ``phase invariance'' concept.
}

\begin{center}
\begin{tabular}{lp{10.5cm}}
\toprule
\textbf{Tier} & \textbf{How ``phase invariance'' is engaged}\\
\midrule
HS     & ``Why doesn't the arrow's overall rotation matter for length?'' ---
         intuitive analogy, no algebra required\\
BegUG  & Compute $|e^{i\theta}c_n|^2 = |c_n|^2$; verify numerically for
         a 2-component example\\
AdvUG  & Prove phase invariance of all probabilities for an arbitrary orthonormal
         basis expansion\\
MSc    & Interpret as well-definedness of the map $|\psi\rangle\mapsto$ probabilities
         on rays in projective space $\mathbb{P}(\mathcal{H})$\\
PhD    & Formalise rays as 1D subspaces; identify the $U(1)$ fiber structure of the
         Hopf fibration $S^1\to S^{2n+1}\to\mathbb{CP}^n$\\
\bottomrule
\end{tabular}
\end{center}

%-------------------------------------------------------------
\subsection{Worked Example: Bra-Ket Module I.2, Lecture L01}

The following shows how a single lecture --- \textbf{Module~I.2, Lecture~L01:
\textit{Why Dirac Notation? States as Rays \& Superposition}} --- is differentiated
across all five tiers. This is the template for every lecture in the programme.

\begin{tcolorbox}[goldbox]
\textbf{Lecture L01 topic:} States, kets, rays, global phase, Born rule preview.
Dual Sakurai (operational) and Dirac (abstract) perspectives.
\end{tcolorbox}

\medskip
\begin{tcolorbox}[tierbox=tier1col, title={Tier 1 (HS) --- L01 engagement}]
\textbf{Key question:} What is a quantum state, informally?

\medskip
\textbf{Content focus:}
\begin{itemize}[leftmargin=1.5em,itemsep=2pt]
  \item $|z|^2$ as probability; interference example with two paths
  \item Ket as a ``direction arrow'' in an abstract space
  \item Born rule as ``probability = squared amplitude''
  \item Why rotating the arrow globally (multiplying by $e^{i\theta}$) does
        not change any measurement outcome
\end{itemize}

\textbf{Signature problem:}
Given $A_1=\tfrac{1}{\sqrt{2}}$, $A_2=\tfrac{i}{\sqrt{2}}$, compute the
total probability for a detector that sums both amplitudes. Verify it equals 1.
\end{tcolorbox}

\medskip
\begin{tcolorbox}[tierbox=tier2col, title={Tier 2 (BegUG) --- L01 engagement}]
\textbf{Key question:} How do I write and work with quantum states?

\medskip
\textbf{Content focus:}
\begin{itemize}[leftmargin=1.5em,itemsep=2pt]
  \item Normalisation $\langle\psi|\psi\rangle=1$; components $c_n=\langle n|\psi\rangle$
  \item Notation dictionary: ket $\leftrightarrow$ column vector $\leftrightarrow$ wavefunction
  \item Build a 2-state example (spin-$\tfrac{1}{2}$) explicitly
\end{itemize}

\textbf{Signature problem:}
Write $|\psi\rangle = \tfrac{3}{5}|+\rangle+\tfrac{4}{5}|-\rangle$.
Verify normalization. Find $P(+)$ and $P(-)$. Change the global phase by
$e^{i\pi/3}$ and confirm probabilities are unchanged.
\end{tcolorbox}

\medskip
\begin{tcolorbox}[tierbox=tier3col, title={Tier 3 (AdvUG) --- L01 engagement}]
\textbf{Key question:} Why is global phase exactly unobservable?

\medskip
\textbf{Content focus:}
\begin{itemize}[leftmargin=1.5em,itemsep=2pt]
  \item Cauchy--Schwarz as bound on all overlaps
  \item Prove that $|\psi\rangle$ and $e^{i\theta}|\psi\rangle$ give identical
        probabilities for \emph{any} observable in \emph{any} orthonormal basis
  \item Discuss why \emph{relative} phase (between components) is physical
\end{itemize}

\textbf{Signature problem:}
Prove that for any normalised $|\psi\rangle$, any orthonormal basis $\{|n\rangle\}$,
and any $\theta\in\mathbb{R}$: $|\langle n|e^{i\theta}|\psi\rangle|^2 = |\langle n|\psi\rangle|^2$
for all $n$, and conclude the total probability is also invariant.
\end{tcolorbox}

\medskip
\begin{tcolorbox}[tierbox=tier4col, title={Tier 4 (MSc) --- L01 engagement}]
\textbf{Key question:} What is the correct mathematical object for a quantum state?

\medskip
\textbf{Content focus:}
\begin{itemize}[leftmargin=1.5em,itemsep=2pt]
  \item Rays as equivalence classes; projective Hilbert space $\mathbb{P}(\mathcal{H})$
  \item The Born-rule map $\mathbb{P}(\mathcal{H})\to[0,1]$ is well-defined on rays
  \item Preview: Wigner's theorem classifies symmetries of $\mathbb{P}(\mathcal{H})$
\end{itemize}

\textbf{Signature problem:}
Define the equivalence relation on $\mathcal{H}\setminus\{0\}$ that defines rays.
Show the Born-rule probability $|\langle a|\psi\rangle|^2$ depends only on the ray $[\psi]$,
not the representative $|\psi\rangle$. What would it mean for a map between physical states
to be a ``symmetry''?
\end{tcolorbox}

\medskip
\begin{tcolorbox}[tierbox=tier5col, title={Tier 5 (PhD) --- L01 engagement}]
\textbf{Key question:} What foundational structures does quantum mechanics actually require?

\medskip
\textbf{Content focus:}
\begin{itemize}[leftmargin=1.5em,itemsep=2pt]
  \item Precise Hilbert space definition; role of completeness and separability
  \item Rays as 1D subspaces; $\mathbb{P}(\mathcal{H})$ as a metric space
        (Fubini--Study metric)
  \item Wigner's theorem (statement in full): every symmetry of $\mathbb{P}(\mathcal{H})$
        lifts to a unitary or antiunitary operator on $\mathcal{H}$
\end{itemize}

\textbf{Signature problem:}
(a) State the definition of a separable Hilbert space.
(b) Why does separability guarantee a countable orthonormal basis?
(c) State Wigner's theorem precisely, including the lifting ambiguity ($\pm$ and $U(1)$ phase).
(d) Give one example of an antiunitary symmetry in QM and explain its physical significance.
\end{tcolorbox}

\newpage
%#############################################################
\section{Assessment Artifact Set}
\label{sec:artifacts}
%#############################################################

\begin{tcolorbox}[goldbox]
\textbf{Design goals.} Each lecture should assess \emph{four distinct capabilities}:
\begin{enumerate}[leftmargin=2em,itemsep=2pt]
  \item \textbf{Conceptual understanding} --- can you explain it?
  \item \textbf{Procedural fluency} --- can you compute it?
  \item \textbf{Structural mastery} --- can you prove or derive it?
  \item \textbf{Transfer \& creation} --- can you apply it to a new context, or
        ask new questions?
\end{enumerate}

\medskip
The six artifact types defined in this section map cleanly onto these capabilities and
are reusable across every lecture of every module in the programme.
\end{tcolorbox}

%-------------------------------------------------------------
\subsection{Capability--Artifact Mapping}

\begin{center}
\begin{tabularx}{\textwidth}{lX}
\toprule
\textbf{Capability} & \textbf{Artifact(s)}\\
\midrule
Conceptual understanding  & Concept Questions (CQ)\\
Procedural fluency        & Simple Problem Set (SPS)\\
Structural mastery        & Complex Problem Set (CPS)\\
Transfer (engineering)    & Simple Project (SPJ)\\
Transfer (synthesis)      & Complex Project (CPJ)\\
Creation \& frontier thinking & Well-Defined Research Question (WRQ);
                                Open-Ended Research Question (ORQ)\\
\bottomrule
\end{tabularx}
\end{center}

%-------------------------------------------------------------
\subsection{Artifact 1: Concept Questions (CQ)}

\begin{tcolorbox}[artifactbox, title={CQ --- Concept Questions}]

\textbf{Purpose:} Fast checks that a student has the \emph{right mental model},
not just algebra.

\medskip
\textbf{Deliverable format:}
\begin{itemize}[leftmargin=1.5em,itemsep=2pt]
  \item 6--10 short questions: mix of multiple-choice, true/false + justify,
        and 2--3 sentence prompts
  \item Designed for 10--15 minutes in class or as pre-lecture warmup
\end{itemize}

\medskip
\textbf{Recommended counts by tier:}
\begin{center}
\begin{tabular}{ll}
\toprule
\textbf{Tier} & \textbf{Question count \& character}\\
\midrule
HS    & 6 (mostly qualitative)\\
BegUG & 8 (qualitative + very short computations)\\
AdvUG & 10 (include ``explain why'')\\
MSc / PhD & 10 (include ``state assumptions'' + ``what fails if\ldots'')\\
\bottomrule
\end{tabular}
\end{center}
\end{tcolorbox}

\medskip
\begin{tcolorbox}[rubricbox, title={CQ Rubric}]
\begin{itemize}[leftmargin=1.5em,itemsep=2pt]
  \item \textbf{2 points each:} correct answer \emph{and} explanation consistent
  \item Total: 12--20 points depending on count
\end{itemize}
\end{tcolorbox}

\medskip
\instructornote{%
  Example CQ prompts (Module~I.2 Lecture~L01 style):\\[2pt]
  (1)~Why is $|\psi\rangle$ not ``a column vector'' until a basis is chosen?\\
  (2)~What is physically meaningful: $\langle a|\psi\rangle$ or $|\langle a|\psi\rangle|^2$?\\
  (3)~Why does a global phase not change probabilities?\\
  (4)~In what sense is ``superposition'' basis-dependent?
}

%-------------------------------------------------------------
\subsection{Artifact 2: Simple Problem Set (SPS)}

\begin{tcolorbox}[artifactbox, title={SPS --- Simple Problem Set}]

\textbf{Purpose:} Build \emph{fluency}: compute probabilities, normalise states,
insert identity, perform basic basis changes.

\medskip
\textbf{Deliverable format:}
\begin{itemize}[leftmargin=1.5em,itemsep=2pt]
  \item 8--12 problems
  \item Mostly 1--5 steps each
  \item Minimal proofs; focus on correct use of notation and technique
\end{itemize}

\medskip
\textbf{Suggested internal structure:}
\begin{itemize}[leftmargin=1.5em,itemsep=2pt]
  \item \textbf{Part A --- mechanics:} normalisation, components, overlaps
  \item \textbf{Part B --- interpretation:} what does the computed number mean?
  \item \textbf{Part C --- translation:}
        ket $\leftrightarrow$ coordinates $\leftrightarrow$ wavefunction preview
\end{itemize}

\medskip
\textbf{Tier scaling:}
\begin{center}
\begin{tabular}{lp{9cm}}
\toprule
HS    & Numeric complex arithmetic + simple amplitude-to-probability\\
BegUG & Normalisation + components $c_n=\langle n|\psi\rangle$ + probabilities\\
AdvUG & Include ``show that'' using completeness, but not full proofs\\
MSc/PhD & Fewer but deeper; emphasise representation-independence checks\\
\bottomrule
\end{tabular}
\end{center}
\end{tcolorbox}

\medskip
\begin{tcolorbox}[rubricbox, title={SPS Rubric}]
\begin{tabular}{ll}
  Correct result         & 60\%\\
  Correct method/notation & 30\%\\
  Interpretation sentence & 10\%\\
\end{tabular}
\end{tcolorbox}

%-------------------------------------------------------------
\subsection{Artifact 3: Complex Problem Set (CPS)}

\begin{tcolorbox}[artifactbox, title={CPS --- Complex Problem Set}]

\textbf{Purpose:} Assess \emph{structural mastery}: proofs, derivations, multi-step
synthesis, ``why the formalism works.''

\medskip
\textbf{Deliverable format:}
\begin{itemize}[leftmargin=1.5em,itemsep=2pt]
  \item 4--6 problems, multi-part (a)(b)(c)\ldots
  \item Often proof-based or derivation-based
  \item Designed for 3--6 hours total work
\end{itemize}

\medskip
\textbf{Canonical CPS themes in early lectures:}
\begin{itemize}[leftmargin=1.5em,itemsep=2pt]
  \item Prove Cauchy--Schwarz and/or triangle inequality
  \item Prove phase invariance and ray equivalence
  \item Show unitaries preserve probabilities and inner products
  \item Derive completeness consequences (Parseval identity)
\end{itemize}

\medskip
\textbf{Tier scaling:}
\begin{center}
\begin{tabular}{lp{9cm}}
\toprule
AdvUG & Full proofs with guidance allowed\\
MSc   & Proofs + conceptual interpretation + links to projective space\\
PhD   & Proofs + additional conditions (boundedness, domains, separability)\\
\bottomrule
\end{tabular}
\end{center}
\end{tcolorbox}

\medskip
\begin{tcolorbox}[rubricbox, title={CPS Rubric}]
\begin{tabular}{ll}
  Correct statement of assumptions & 15\%\\
  Logical proof structure           & 45\%\\
  Correct algebra/steps             & 25\%\\
  Insight / interpretation          & 15\%\\
\end{tabular}
\end{tcolorbox}

%-------------------------------------------------------------
\subsection{Artifact 4: Simple Project (SPJ)}

\begin{tcolorbox}[artifactbox, title={SPJ --- Simple Project}]

\textbf{Purpose:} Hands-on transfer: ``I can build something that uses the formalism.''

\medskip
\textbf{Scope:} 2--6 hours; structured, well-scaffolded; works for mixed cohorts.

\medskip
\textbf{Deliverable format:}
\begin{itemize}[leftmargin=1.5em,itemsep=2pt]
  \item Short report (1--3 pages) \emph{and} a small artifact:
        code, worksheet, or visualisation
\end{itemize}
\end{tcolorbox}

\medskip
\begin{tcolorbox}[goldbox]
\textbf{Example SPJ Template --- Basis Expansion Calculator (Module~I.2, L01)}

\medskip
\noindent\textbf{Inputs:} coefficients $c$ in basis $\{|n\rangle\}$; unitary matrix $U$.\\
\textbf{Outputs:} new coefficients $c'=Uc$.\\[4pt]
\textbf{Required checks:}
\begin{itemize}[leftmargin=1.5em,itemsep=2pt]
  \item Normalisation preserved: $\|c'\|^2=\|c\|^2=1$
  \item Probabilities sum to 1 in new basis
  \item Global phase invariance: multiplying $c$ by $e^{i\theta}$ leaves all $|c'_n|^2$
        unchanged
\end{itemize}
\textbf{Deliverables:} code + README; one worked example; short explanation of what
changed and what did not.
\end{tcolorbox}

\medskip
\begin{tcolorbox}[rubricbox, title={SPJ Rubric}]
\begin{tabular}{ll}
  Correctness                    & 40\%\\
  Verification tests included    & 25\%\\
  Clarity of explanation         & 25\%\\
  Reproducibility (instructions) & 10\%\\
\end{tabular}
\end{tcolorbox}

%-------------------------------------------------------------
\subsection{Artifact 5: Complex Project (CPJ)}

\begin{tcolorbox}[artifactbox, title={CPJ --- Complex Project}]

\textbf{Purpose:} Synthesis and conceptual transfer: ``I can compare frameworks and
defend choices.''

\medskip
\textbf{Scope:} 8--15 hours; open-ended but guided; should produce a mini-essay or
mini-tool with reflection.

\medskip
\textbf{Deliverable format:}
\begin{itemize}[leftmargin=1.5em,itemsep=2pt]
  \item 5--8 page report + appendix with computations + optional code for
        basis transforms/plots
\end{itemize}
\end{tcolorbox}

\medskip
\begin{tcolorbox}[goldbox]
\textbf{Example CPJ Template --- Compare Representations (Module~I.2, L01)}

\medskip
\noindent Choose a state and represent it in four ways:
\begin{enumerate}[leftmargin=2em,itemsep=2pt]
  \item Abstract ket $|\psi\rangle$ (representation-free)
  \item Coordinates in basis A: $c_n=\langle n|\psi\rangle$
  \item Coordinates in basis B (unitarily related to A)
  \item Preview: wavefunction $\psi(x)=\langle x|\psi\rangle$ (conceptual only)
\end{enumerate}

\noindent\textbf{Analyse:}
\begin{itemize}[leftmargin=1.5em,itemsep=2pt]
  \item What is \emph{invariant} (norm, transition probabilities)?
  \item What \emph{changes} (coordinates, basis-dependent ``superposition'')?
\end{itemize}
\end{tcolorbox}

\medskip
\begin{tcolorbox}[rubricbox, title={CPJ Rubric}]
\begin{tabular}{ll}
  Technical correctness    & 30\%\\
  Conceptual synthesis     & 35\%\\
  Depth \& structure       & 25\%\\
  Originality / insight    & 10\%\\
\end{tabular}
\end{tcolorbox}

%-------------------------------------------------------------
\subsection{Artifact 6: Well-Defined Research Question (WRQ)}

\begin{tcolorbox}[artifactbox, title={WRQ --- Well-Defined Research Question}]

\textbf{Purpose:} Train ``research-like thinking'' while still being \emph{bounded}:
clear scope, clear output, answerable within a few days.

\medskip
\textbf{Scope:} 5--10 hours; must be answerable with literature + calculation +
interpretation.

\medskip
\textbf{Deliverable format --- 2--6 page memo containing:}
\begin{itemize}[leftmargin=1.5em,itemsep=2pt]
  \item Question statement
  \item Method (how you will answer it)
  \item Evidence (citations + computations)
  \item Conclusion
  \item 5--10 references
\end{itemize}

\medskip
\textbf{Example prompts (Module~I.2, early bra-ket):}
\begin{itemize}[leftmargin=1.5em,itemsep=2pt]
  \item ``How did Born's probabilistic interpretation change the meaning of the
        wavefunction between 1926--1927?''
  \item ``What exactly does `state is a ray' buy us operationally?
        Provide examples where vectors fail but rays succeed.''
\end{itemize}
\end{tcolorbox}

\medskip
\begin{tcolorbox}[rubricbox, title={WRQ Rubric}]
\begin{tabular}{ll}
  Clear scope and method    & 25\%\\
  Evidence quality          & 25\%\\
  Correct technical content & 25\%\\
  Quality of argument       & 25\%\\
\end{tabular}
\end{tcolorbox}

%-------------------------------------------------------------
\subsection{Artifact 7: Open-Ended Research Question (ORQ)}

\begin{tcolorbox}[artifactbox, title={ORQ --- Open-Ended Research Question}]

\textbf{Purpose:} Expose the frontier: no single correct answer; evaluates
\emph{quality of reasoning}, not correctness of conclusion.

\medskip
\textbf{Scope:} 10--20 hours (or ongoing); may include proposing a new definition,
comparing frameworks, or designing a thought experiment.

\medskip
\textbf{Deliverable format --- 4--10 page proposal-style note:}
\begin{itemize}[leftmargin=1.5em,itemsep=2pt]
  \item Motivation
  \item What is known (brief review)
  \item What you propose
  \item How you would test or argue it
  \item Risks and limitations
\end{itemize}

\medskip
\textbf{Example prompts (Module~I.2, early bra-ket):}
\begin{itemize}[leftmargin=1.5em,itemsep=2pt]
  \item ``What minimal operational axioms force complex Hilbert space rather than real?''
  \item ``Can you formulate measurement and state update without projective collapse
        but still reproduce observed statistics?''
\end{itemize}
\end{tcolorbox}

\medskip
\begin{tcolorbox}[rubricbox, title={ORQ Rubric}]
\begin{tabular}{ll}
  Clarity + framing             & 25\%\\
  Depth of engagement           & 25\%\\
  Novelty and plausibility      & 25\%\\
  Awareness of limitations      & 25\%\\
\end{tabular}
\end{tcolorbox}

%-------------------------------------------------------------
\subsection{Exercise Packaging Standard (Multi-page Format)}

\instructornote{%
  Every individual problem --- whether in an SPS or CPS --- should be packaged in
  the following four-layer format. This supports self-study, allows automated
  assembly into different deliverables, and makes solution-writing consistent.
}

\begin{tcolorbox}[goldbox]
\textbf{Four-layer exercise packet structure} (one file-set per problem):

\medskip
\begin{tabular}{lp{10cm}}
\toprule
\textbf{Layer} & \textbf{Content}\\
\midrule
1. Problem Statement   & Clean, no hints; the only document students receive at the
                         start\\
2. Suggestions / Hints & Scaffolded prompts; released at student request or after
                         first attempt\\
3. Advanced Suggestions & Structural hints, alternative approaches, connections to
                          other results\\
4. Detailed Solution   & Full derivation + commentary + common pitfall warnings\\
\bottomrule
\end{tabular}

\medskip
\noindent\textbf{File naming convention for repository automation:}\\
\texttt{ex\#\#\_problem.tex}, \texttt{ex\#\#\_hints.tex},
\texttt{ex\#\#\_adv\_hints.tex}, \texttt{ex\#\#\_solution.tex}
\end{tcolorbox}

%-------------------------------------------------------------
\subsection{Recommended Per-Lecture Minimum Bundle}

\begin{tcolorbox}[goldbox]
For each 90-minute lecture, a sustainable and scalable package is:

\medskip
\begin{center}
\begin{tabular}{lr}
\toprule
\textbf{Artifact} & \textbf{Count}\\
\midrule
Concept Questions (CQ)           & 8\\
Simple Problem Set (SPS)         & 10 problems\\
Complex Problem Set (CPS)        & 4 problems\\
Simple Project (SPJ)             & 1\\
Complex Project (CPJ)            & 1\\
Well-Defined Research Q.\ (WRQ)  & 1\\
Open-Ended Research Q.\ (ORQ)    & 1\\
\bottomrule
\end{tabular}
\end{center}

\medskip
\noindent This yields:
\begin{itemize}[leftmargin=1.5em,itemsep=2pt]
  \item Enough material for a mixed cohort (simple through PhD tracks)
  \item A manageable grading load per lecture cycle
  \item A consistent cadence across the entire programme
\end{itemize}
\end{tcolorbox}

%-------------------------------------------------------------
\subsection{Suggested Assessment Weights}

\begin{center}
\begin{tabular}{lrp{7cm}}
\toprule
\textbf{Artifact type} & \textbf{Weight} & \textbf{Rationale}\\
\midrule
Concept Questions (CQ) & 10\% & Fast feedback; keeps students on track\\
Simple Problem Set (SPS) & 30\% & Core fluency; highest volume\\
Complex Problem Set (CPS) & 30\% & Structural mastery; differentiates tracks\\
Projects (SPJ + CPJ) & 20\% & Transfer and synthesis; project-based learning\\
Research Questions (WRQ + ORQ) & 10\% & Frontier thinking; rewards initiative\\
\midrule
\textbf{Total} & \textbf{100\%} & \\
\bottomrule
\end{tabular}
\end{center}

\medskip
\noindent\textbf{Track-based grading policy:}
\begin{itemize}[leftmargin=1.5em]
  \item HS / BegUG $\to$ \textbf{Simple track}: CQ + SPS + SPJ only
  \item AdvUG / MSc / PhD $\to$ \textbf{Complex track}: all artifacts;
        CPS and CPJ are mandatory
\end{itemize}

\newpage
%=============================================================
\section*{Appendix --- Quick Reference Card}
\addcontentsline{toc}{section}{Appendix --- Quick Reference Card}
%=============================================================

\begin{center}
\textbf{\large Tier $\times$ Artifact Grid}\\[6pt]
\small
\begin{tabular}{l|ccccccc}
\toprule
 & \textbf{CQ} & \textbf{SPS} & \textbf{CPS} & \textbf{SPJ} & \textbf{CPJ} & \textbf{WRQ} & \textbf{ORQ}\\
\midrule
HS    & \checkmark & \checkmark &            & \checkmark &            &            &\\
BegUG & \checkmark & \checkmark &            & \checkmark & (optional) &            &\\
AdvUG & \checkmark & \checkmark & \checkmark & \checkmark & \checkmark & \checkmark &\\
MSc   & \checkmark & \checkmark & \checkmark & \checkmark & \checkmark & \checkmark & \checkmark\\
PhD   & \checkmark & \checkmark & \checkmark & \checkmark & \checkmark & \checkmark & \checkmark\\
\bottomrule
\end{tabular}
\end{center}

\bigskip
\begin{center}
\textbf{\large Artifact Duration Summary}\\[6pt]
\begin{tabular}{llll}
\toprule
\textbf{Artifact} & \textbf{Format} & \textbf{Est.\ time} & \textbf{Primary output}\\
\midrule
CQ  & 6--10 questions & 10--15 min   & Answers + brief justifications\\
SPS & 8--12 problems  & 2--4 hours   & Worked solutions\\
CPS & 4--6 problems   & 3--6 hours   & Proofs + derivations\\
SPJ & Structured task & 2--6 hours   & Code / worksheet + 1--3 pg report\\
CPJ & Open task       & 8--15 hours  & 5--8 pg report + appendix\\
WRQ & Scoped question & 5--10 hours  & 2--6 pg memo\\
ORQ & Open question   & 10--20 hours & 4--10 pg proposal\\
\bottomrule
\end{tabular}
\end{center}

%=============================================================
\end{document}
